% ═══════════════════════════════════════════════════════════════════════════════
% CAPÍTULO 4 — REDES NEURONALES: EL APROXIMADOR UNIVERSAL
% ═══════════════════════════════════════════════════════════════════════════════
% Cambios principales:
%   1. Apertura mostrando que las familias fijas fallan en alta dimensión y
%      que las redes neuronales son familias con bases adaptativas.
%   2. Activaciones con física integrada: sigmoide = Fermi-Dirac, tanh =
%      función de Langevin paramagnética. No como recuadros aislados sino
%      como parte de la derivación.
%   3. Teorema de aproximación universal con discusión de profundidad vs.
%      anchura (multiescala como analogía).
%   4. Oscilador: MLP vs. métodos lineales, cierre del arco de 4 capítulos.
%   5. Transición fuerte al Cap. 5: "tenemos la familia; ahora necesitamos
%      el algoritmo de optimización y su física".
%   6. Estrella Polar independiente al final.
% ═══════════════════════════════════════════════════════════════════════════════

\chapter{Redes neuronales: el aproximador universal}
\label{ch:cap4}

% ─────────────────────────────────────────────────────────────────────────────
\section{De las bases fijas a las bases adaptativas}
\label{sec:bases_adaptativas}
% ─────────────────────────────────────────────────────────────────────────────

\dificultad{1}{2}

Los tres capítulos anteriores construyeron las herramientas del ajuste
clásico: familias paramétricas (Cap.~\ref{ch:cap1}), control del
sobreajuste (Cap.~\ref{ch:cap2}), y geometría del paisaje de pérdida
(Cap.~\ref{ch:cap3}). En todos estos capítulos, las familias lineales
—polinomios, series de Fourier, funciones radiales— resolvieron el
problema del oscilador armónico de forma satisfactoria. Pero el
oscilador es un problema unidimensional con estructura conocida.

Volvamos al problema del potencial interatómico de la
Sección~\ref{sec:motivacion1}: una función de $\mathbb{R}^{192}$ que
depende de las posiciones de 64 átomos. Una expansión en monomios de
grado $d$ requiere $\binom{192+d}{d}$ funciones base. Para $d = 3$:
$\sim 10^6$ funciones. Para $d = 5$: $\sim 10^{10}$. El número de
funciones base crece polinomialmente con la dimensión $n$ para grado
fijo, y la constante es astronómica. Ningún conjunto de datos finito
puede determinar $10^{10}$ coeficientes de forma útil.

La maldición de la dimensionalidad no es un problema de computación
sino de \emph{representación}: las bases fijas son ineficientes porque
asignan la misma cantidad de recursos a todas las regiones del espacio
de entrada, incluyendo las regiones vacías de datos. La solución es una
familia que concentre sus ``funciones base'' donde los datos están, y
que esas funciones base se \emph{aprendan} de los datos.

Una red neuronal es exactamente eso: una familia paramétrica donde las
funciones base son \textbf{adaptativas}. En su forma más simple, un
\textbf{perceptrón multicapa} (MLP) de una capa oculta con $m$ neuronas
es:
\begin{equation}
  f(\bm{x};\bm{\theta})
  = \sum_{j=1}^m c_j\,\sigma(\bm{w}_j^\top\bm{x} + b_j),
  \label{eq:mlp_una_capa}
\end{equation}
donde $\sigma$ es una función no lineal (activación), $\bm{w}_j \in
\mathbb{R}^n$ son los pesos, $b_j$ los sesgos y $c_j$ los coeficientes
de la capa de salida. Los parámetros
$\bm{\theta} = \{c_j, \bm{w}_j, b_j\}_{j=1}^m$ determinan
simultáneamente la \emph{forma} de las funciones base
$\phi_j(\bm{x}) = \sigma(\bm{w}_j^\top\bm{x}+b_j)$ y sus
\emph{coeficientes} $c_j$. Esta doble adaptabilidad es la clave de su
eficiencia.

% ─────────────────────────────────────────────────────────────────────────────
\section{La neurona artificial y las funciones de activación}
\label{sec:neurona}
% ─────────────────────────────────────────────────────────────────────────────

\dificultad{2}{2}

La unidad básica de una red neuronal es la \textbf{neurona artificial}:
\begin{equation}
  a = \sigma(z), \qquad z = \bm{w}^\top\bm{x} + b,
  \label{eq:neurona}
\end{equation}
donde $z$ es la \textbf{preactivación} (combinación lineal de las
entradas) y $a$ es la \textbf{activación} (salida después de la no
linealidad). La función $\sigma$ introduce la no linealidad que
distingue a la red de una familia lineal.

\subsection{Funciones de activación y sus conexiones con la física}
\label{subsec:activaciones}

Las funciones de activación más usadas tienen interpretaciones físicas
directas que ayudan a entender su comportamiento.

\textbf{La sigmoide logística:}
\begin{equation}
  \sigma(z) = \frac{1}{1+e^{-z}}.
  \label{eq:sigmoide}
\end{equation}
Un físico de estado sólido reconoce esta función inmediatamente: es la
\textbf{distribución de Fermi-Dirac} $n_F(\varepsilon) = 1/(1 +
e^{(\varepsilon-\mu)/(k_BT)})$ con $\mu = -b/w$ (potencial químico) y
$k_BT = 1/w$ (temperatura). A temperatura cero, la función es un escalón
(Heaviside): todos los estados por debajo de $\mu$ están ocupados, todos
los de arriba están vacíos. A temperatura finita, la transición se suaviza
en una región de ancho $\sim k_BT$ alrededor de $\mu$. En la neurona, el
peso $w$ controla la ``nitidez'' de la transición y el sesgo $b$ controla
su posición. Una red de neuronas sigmoide construye funciones arbitrarias
superponiendo escalones suavizados en distintas posiciones y con distintas
nitideces: la misma idea que la reconstrucción de una distribución de
energía a partir de la curva de Fermi.

\textbf{La tangente hiperbólica:}
\begin{equation}
  \tanh(z) = \frac{e^z - e^{-z}}{e^z + e^{-z}} = 2\sigma(2z) - 1.
  \label{eq:tanh}
\end{equation}
Es la \textbf{función de Langevin del paramagnetismo clásico}: la
magnetización de un sistema de espines clásicos en un campo externo $h$
es $\langle m \rangle = \tanh(\beta h)$, donde $\beta = 1/(k_BT)$. A
diferencia de la sigmoide (que mapea a $[0,1]$), la tanh mapea a
$[-1,1]$ y tiene media cero para entradas simétricas, lo que facilita
la propagación del gradiente.

\textbf{La ReLU (Rectified Linear Unit):}
\begin{equation}
  \text{ReLU}(z) = \max(0, z).
  \label{eq:relu}
\end{equation}
La ReLU no tiene una analogía termodinámica directa, pero tiene una
interpretación geométrica poderosa: cada neurona ReLU divide el
espacio de entrada en dos regiones (activa e inactiva) separadas por
un hiperplano $\bm{w}^\top\bm{x}+b = 0$. Una red de $m$ neuronas
ReLU en una capa crea hasta $\binom{m}{n}$ regiones lineales por
trozos, cada una con su propia función afín. La red construye
funciones complejas pegando regiones lineales, como un mosaico.

La ReLU tiene una ventaja computacional decisiva: su derivada es $0$ o
$1$ (no exponenciales), lo que previene el \textbf{problema del
gradiente que se desvanece} (\textit{vanishing gradient}). Con
activaciones sigmoide o tanh, la derivada $\sigma'(z) \in (0, 0.25]$
se multiplica capa tras capa: en una red de $L$ capas, el gradiente
puede escalar como $(0.25)^L$, que para $L = 20$ es $\sim 10^{-12}$.
Con ReLU, las neuronas activas transmiten el gradiente sin atenuación:
$\text{ReLU}'(z) = 1$ para $z > 0$. Esta propiedad hizo posible
entrenar redes profundas ($L \gg 10$) por primera vez.

% ─────────────────────────────────────────────────────────────────────────────
\section{El perceptrón multicapa (MLP)}
\label{sec:mlp}
% ─────────────────────────────────────────────────────────────────────────────

\dificultad{2}{2}

\begin{definicion}{Perceptrón multicapa (MLP)}
\label{def:mlp}
Un MLP de $L$ capas con anchuras
$(n_0, n_1, \ldots, n_L)$ es la composición:
\begin{align}
  \bm{z}^{(l)} &= \bm{W}^{(l)}\bm{a}^{(l-1)} + \bm{b}^{(l)},
  \quad l = 1, \ldots, L, \label{eq:preact}\\
  \bm{a}^{(l)} &= \sigma(\bm{z}^{(l)}),
  \quad l = 1, \ldots, L-1, \label{eq:act}\\
  f(\bm{x};\bm{\theta}) &= \bm{z}^{(L)},
  \label{eq:salida}
\end{align}
con $\bm{a}^{(0)} = \bm{x}$, $\bm{W}^{(l)} \in
\mathbb{R}^{n_l \times n_{l-1}}$, $\bm{b}^{(l)} \in \mathbb{R}^{n_l}$.
La última capa es lineal (sin activación) para regresión.
\end{definicion}

La notación $(n_0, n_1, \ldots, n_L)$ especifica la arquitectura
completamente. Por ejemplo, $(2, 16, 16, 1)$ es una red con 2 entradas,
dos capas ocultas de 16 neuronas cada una, y una salida escalar. El
número total de parámetros es:
\begin{equation}
  p = \sum_{l=1}^L (n_{l-1}+1)n_l
  = \sum_{l=1}^L n_{l-1}n_l + n_l.
  \label{eq:conteo}
\end{equation}
Para la arquitectura $(2, 16, 16, 1)$: $p = (2+1)\cdot 16 +
(16+1)\cdot 16 + (16+1)\cdot 1 = 48 + 272 + 17 = 337$ parámetros.

La estructura composicional del MLP tiene una interpretación como
sistema dinámico discreto que anticipa las Neural ODEs del
Capítulo~\ref{ch:cap9}. Si las capas tienen la misma anchura
($n_l = m$ para todo $l$) y se añaden conexiones residuales
$\bm{a}^{(l)} = \bm{a}^{(l-1)} + \sigma(\bm{W}^{(l)}\bm{a}^{(l-1)}
+ \bm{b}^{(l)})$, cada capa es un paso del integrador de Euler para la
ODE $\dot{\bm{a}} = F(\bm{a})$. La profundidad $L$ es el número de
pasos de Euler y la red profunda es la discretización de un sistema
dinámico continuo. Esta conexión se desarrolla completamente en el
Capítulo~\ref{ch:cap9}.

% ─────────────────────────────────────────────────────────────────────────────
\section{El teorema de aproximación universal}
\label{sec:universal}
% ─────────────────────────────────────────────────────────────────────────────

\dificultad{3}{2}

\begin{teorema}{Aproximación universal (Cybenko 1989, Hornik 1989)}
\label{thm:universal}
Sea $\sigma: \mathbb{R} \to \mathbb{R}$ continua y no constante. Para
toda función continua $g: K \to \mathbb{R}$ definida en un compacto
$K \subset \mathbb{R}^n$ y para todo $\varepsilon > 0$, existe un
entero $m$ y parámetros $\{c_j, \bm{w}_j, b_j\}_{j=1}^m$ tales que:
\begin{equation}
  \sup_{\bm{x} \in K}\left|g(\bm{x})
  - \sum_{j=1}^m c_j\sigma(\bm{w}_j^\top\bm{x}+b_j)\right|
  < \varepsilon.
  \label{eq:universal}
\end{equation}
\end{teorema}

El teorema dice que una red de una sola capa oculta puede aproximar
cualquier función continua con precisión arbitraria, siempre que se
usen suficientes neuronas $m$. Pero no dice cuántas neuronas se
necesitan. Y ahí está la trampa: para funciones generales en
$\mathbb{R}^n$, el número $m$ necesario puede crecer exponencialmente
con $n$. El teorema garantiza la existencia de una buena aproximación
pero no su eficiencia.

\subsection{Profundidad vs.\ anchura}
\label{subsec:profundidad}

La ventaja de la profundidad (más capas) sobre la anchura (más neuronas
por capa) es análoga a la ventaja de los métodos multiescala sobre los
de una sola escala en física computacional. Un problema con estructura
jerárquica —por ejemplo, la distribución de carga en un cristal, que
tiene estructura a escalas atómica, de celda unitaria y macroscópica—
se resuelve eficientemente con métodos que operan simultáneamente a
múltiples escalas (multigrid, wavelets, renormalización). Una red
profunda hace algo similar: las capas iniciales detectan características
de baja escala (bordes, frecuencias locales), y las capas posteriores
combinan esas características en patrones de escala mayor.

\begin{proposicion}{Ventaja de la profundidad}
\label{prop:profundidad}
Existen funciones en $\mathbb{R}^n$ que un MLP de profundidad $L$ con
$O(n)$ neuronas por capa puede representar con $O(nL)$ parámetros,
pero que una red de profundidad 1 necesita $\Omega(2^n)$ neuronas para
representar con la misma precisión.
\end{proposicion}

Un ejemplo clásico es la función paridad en $\{0,1\}^n$: una red de
profundidad $\log n$ con $O(n)$ neuronas por capa la computa
exactamente, pero una red de profundidad 1 necesita $2^{n-1}$ neuronas.

La profundidad también tiene un costo: paisajes de pérdida más
complicados (Capítulo~\ref{ch:cap3}), gradientes que se desvanecen o
explotan, y mayor dificultad de entrenamiento. Las técnicas para
manejar estos problemas —conexiones residuales, normalización,
inicialización cuidadosa— se desarrollan en los
Capítulos~\ref{ch:cap5}--\ref{ch:cap8}.

% ─────────────────────────────────────────────────────────────────────────────
\section{El oscilador como laboratorio: la red neuronal vs.\ los
  métodos lineales}
\label{sec:oscilador4}
% ─────────────────────────────────────────────────────────────────────────────

\dificultad{2}{2}

Completamos el arco del ejemplo unificado. Los datos son los mismos de
los capítulos anteriores: $N = 25$ puntos de $\sin(2\pi t) +
\varepsilon_i$ con $\sigma = 0.15$ en $[0, 2]$.

\textbf{Resumen de los métodos anteriores:}
\begin{center}
\begin{tabular}{llll}
\toprule
\textbf{Método} & \textbf{Cap.} & $\mathcal{L}_{\text{test}}$ &
\textbf{Observación} \\
\midrule
Monomios, $p = 6$ & \ref{ch:cap1} & $\sim 0.04$ & Sesgo moderado \\
Fourier, $K = 3$ & \ref{ch:cap1} & $\sim 0.02$ & Óptima si se
  conoce la estructura \\
Monomios, $p = 20$, sin reg. & \ref{ch:cap2} & $> 10$ & Sobreajuste
  severo \\
Monomios, $p = 20$, $\lambda = 10^{-4}$ & \ref{ch:cap2} & $\sim 0.03$
  & Regularización efectiva \\
\bottomrule
\end{tabular}
\end{center}

\textbf{Con un MLP:} un MLP $(1, 32, 1)$ con activación ReLU tiene
$p = (1+1)\cdot 32 + (32+1)\cdot 1 = 97$ parámetros. Entrenado con el
algoritmo Adam (Capítulo~\ref{ch:cap5}) durante 2000 épocas con tasa de
aprendizaje $\eta = 10^{-3}$, obtiene $\mathcal{L}_{\text{test}} \approx
0.025$: comparable al mejor método lineal (Fourier) pero sin haber
especificado la estructura armónica de la función.

La clave es que la red ``descubre'' las funciones base apropiadas: las
neuronas ReLU se posicionan automáticamente para formar una aproximación
lineal por tramos de la función seno. Si se examina la preactivación
$z_j = w_j t + b_j$ de cada neurona, se observa que los pesos $w_j$ y
sesgos $b_j$ se organizan para que los ``codos'' de la ReLU coincidan
con los puntos de inflexión y los extremos de $\sin(2\pi t)$.

\textbf{¿Por qué usar una red si Fourier funciona?} Para el oscilador
armónico en 1D, la base de Fourier es imbatible porque captura
exactamente la estructura del problema con un solo coeficiente. Pero
en un problema donde la estructura es desconocida o no armónica —el
potencial interatómico del inicio del Capítulo~\ref{ch:cap1}, por
ejemplo—, las bases de Fourier no ayudan y la red neuronal es la
herramienta que escala.

El Ejercicio Computacional~\ref{ejcomp:oscilador4} implementa esta
comparación.

% ─────────────────────────────────────────────────────────────────────────────
\section{Resumen del Capítulo 4 y transición al Arco II}
\label{sec:resumen4}
% ─────────────────────────────────────────────────────────────────────────────

El MLP es una familia paramétrica no lineal con bases adaptativas que
resuelve el problema de escalabilidad de las familias lineales. El
teorema de aproximación universal garantiza que puede representar
cualquier función continua con suficientes neuronas. La profundidad
aporta eficiencia representacional (jerarquía de escalas) a costa de
paisajes de pérdida más complicados.

Pero hay una pregunta que este capítulo no respondió: ¿cómo se
encuentran los parámetros $\bm{\theta}^*$ que minimizan la pérdida?
Para los 337 parámetros del MLP $(2, 16, 16, 1)$, o los $10^6$ de un
modelo realista, las ecuaciones normales no existen (la familia es no
lineal) y la visualización del paisaje (Capítulo~\ref{ch:cap3}) es
imposible en $\mathbb{R}^p$ para $p > 3$. Se necesita un algoritmo
iterativo que descienda por el paisaje usando solo información local
(gradiente, quizá hessiano).

Ese algoritmo es el \textbf{descenso de gradiente estocástico} (SGD),
y su análisis revela la conexión más profunda del libro. El SGD no es
simplemente un algoritmo de optimización: es la discretización de una
\textbf{ecuación de Langevin}, la ecuación del movimiento de una
partícula browniana en un potencial de energía. La distribución
estacionaria de esa dinámica es la \textbf{distribución de Gibbs} de
la mecánica estadística. La tasa de aprendizaje es la temperatura. El
recocido simulado del learning rate schedule es el enfriamiento lento
de un cristal.

El Capítulo~\ref{ch:cap5} desarrolla esta conexión. No es una analogía:
es una identidad matemática.

% ─────────────────────────────────────────────────────────────────────────────
\section*{Ejercicios resueltos}
\addcontentsline{toc}{section}{Ejercicios resueltos}
% ─────────────────────────────────────────────────────────────────────────────

\begin{ejercicio}
\label{ej:forward_manual}
\textbf{(Forward pass manual).}

Sea un MLP $(2, 3, 1)$ con activación ReLU y parámetros:
$\bm{W}^{(1)} = \begin{pmatrix} 1 & -1 \\ 1 & 1 \\ -1 & 0
\end{pmatrix}$,
$\bm{b}^{(1)} = \begin{pmatrix} -2 \\ 0 \\ -2 \end{pmatrix}$,
$\bm{W}^{(2)} = (1, 1, -1)$, $b^{(2)} = 0$.

Calcular la salida para $\bm{x} = (1, 2)^\top$.

\medskip\noindent\textit{Solución.}
$\bm{z}^{(1)} = \bm{W}^{(1)}\bm{x}+\bm{b}^{(1)}
= (1-2-2, 1+2+0, -1+0-2)^\top = (-3, 3, -3)^\top$.

$\bm{a}^{(1)} = \text{ReLU}(\bm{z}^{(1)}) = (0, 3, 0)^\top$.

$z^{(2)} = \bm{W}^{(2)}\bm{a}^{(1)}+b^{(2)} = 0+3+0 = 3$.

$f(\bm{x};\bm{\theta}) = 3$.

Solo la segunda neurona de la capa oculta está activa ($z_2 > 0$); las
otras dos tienen preactivación negativa y la ReLU las apaga. La red
efectivamente se reduce a $f = a_2^{(1)} = z_2^{(1)} = \bm{w}_2^\top
\bm{x}+b_2 = x_1+x_2$ para esta entrada particular. $\square$
\end{ejercicio}

\begin{ejercicio}
\label{ej:conteo_parametros}
\textbf{(Conteo de parámetros y regiones lineales).}

\medskip\noindent\textit{(a)} Calcular el número de parámetros de las
arquitecturas $(1, 8, 1)$, $(1, 32, 1)$ y $(1, 16, 16, 1)$.

\medskip\noindent\textit{Solución.}
$(1, 8, 1)$: $(1+1)\cdot 8 + (8+1)\cdot 1 = 16 + 9 = 25$.
$(1, 32, 1)$: $(1+1)\cdot 32 + (32+1)\cdot 1 = 64 + 33 = 97$.
$(1, 16, 16, 1)$: $(1+1)\cdot 16 + (16+1)\cdot 16 + (16+1)\cdot 1 =
32 + 272 + 17 = 321$.

\medskip\noindent\textit{(b)} Para una entrada 1D con activación ReLU,
¿cuántas regiones lineales puede crear cada arquitectura?

\medskip\noindent\textit{Solución.}
Para $n = 1$: una capa de $m$ neuronas ReLU crea a lo sumo $m + 1$
regiones. Una red $(1, 16, 16, 1)$ puede crear hasta $16 \times 16 + 1
= 257$ regiones lineales (el producto de las anchuras más uno, cota
superior). La red profunda $(1, 16, 16, 1)$ con 321 parámetros puede
crear más regiones que la red ancha $(1, 32, 1)$ con solo 97 parámetros
(33 regiones): la profundidad multiplica la capacidad combinatoria.
$\square$
\end{ejercicio}

% ─────────────────────────────────────────────────────────────────────────────
\section*{Ejercicios propuestos}
\addcontentsline{toc}{section}{Ejercicios propuestos}
% ─────────────────────────────────────────────────────────────────────────────

\begin{ejercicio}
\textbf{(Funciones de activación).}
\begin{enumerate}[label=(\alph*)]
  \item Demostrar que $\sigma(z) = 1/(1+e^{-z})$ satisface
        $\sigma'(z) = \sigma(z)(1-\sigma(z))$ y que $\max_z \sigma'(z)
        = 1/4$ (alcanzado en $z = 0$).
  \item Demostrar que $\tanh(z) = 2\sigma(2z) - 1$ y que
        $\tanh'(z) = 1 - \tanh^2(z)$.
  \item Para ReLU, verificar numéricamente que $\sigma'(z) = 1$ para
        $z > 0$ y $\sigma'(z) = 0$ para $z < 0$ usando diferencias
        finitas con paso $h = 10^{-5}$.
  \item Simular el producto de $L$ derivadas para $L = 1, \ldots, 20$
        y cada activación, con $z_k \sim \mathcal{N}(0, 1)$. Graficar
        $\mathbb{E}[\prod_{k=1}^L|\sigma'(z_k)|]$ en escala log. ¿Cuál
        activación mantiene el gradiente estable?
\end{enumerate}
\end{ejercicio}

\begin{ejercicio}
\textbf{(La sigmoide como distribución de Fermi-Dirac).}
\begin{enumerate}[label=(\alph*)]
  \item Para la neurona $\sigma(wx+b)$ con $w = 5$ y $b = -10$,
        graficar la activación en $x \in [0, 4]$. ¿Dónde está la
        ``energía de Fermi'' $\mu = -b/w$? ¿Cuál es la ``temperatura''
        $k_BT = 1/w$?
  \item Superponer tres neuronas sigmoide con $(w_j, b_j)$ distintos y
        coeficientes $c_j$ elegidos para aproximar la función escalón
        $g(x) = \text{sgn}(x-1)$ en $[0, 2]$. ¿Cuántas neuronas se
        necesitan para error $< 0.05$?
\end{enumerate}
\end{ejercicio}

\begin{ejercicio}
\textbf{(Profundidad vs.\ anchura para una función simple).}
Considere $g(x) = |x|$ en $[-1, 1]$.
\begin{enumerate}[label=(\alph*)]
  \item Mostrar que una neurona ReLU calcula
        $\text{ReLU}(x) = \max(0, x) = (x + |x|)/2$, y que
        $|x| = \text{ReLU}(x) + \text{ReLU}(-x)$.
  \item ¿Cuántas neuronas ReLU en una capa necesita un MLP $(1, m, 1)$
        para representar $g(x) = |x|$ exactamente?
  \item ¿Cuántas neuronas necesita un MLP $(1, m, m, 1)$ de dos capas?
\end{enumerate}
\end{ejercicio}

% ─────────────────────────────────────────────────────────────────────────────
\section*{Ejercicios computacionales}
\addcontentsline{toc}{section}{Ejercicios computacionales}
% ─────────────────────────────────────────────────────────────────────────────

\begin{ejerciciocomp}{El oscilador armónico: MLP vs.\ métodos lineales}
\label{ejcomp:oscilador4}
\textbf{Objetivo:} entrenar un MLP para aproximar la función del
oscilador y comparar con los métodos lineales de los capítulos
anteriores.

\medskip\noindent\textbf{Datos:}
$N_{\text{train}} = 25$ puntos de $\sin(2\pi t) + \varepsilon_i$,
$\sigma = 0.15$, $t_i \in [0, 2]$. $N_{\text{test}} = 200$ puntos
adicionales.

\medskip\noindent\textbf{Algoritmo:}
\begin{enumerate}
  \item Entrenar tres arquitecturas con activación ReLU:
        (A)~$(1, 8, 1)$, (B)~$(1, 32, 1)$, (C)~$(1, 16, 16, 1)$.
        Usar MSE como pérdida y Adam con $\eta = 10^{-3}$ durante
        $2000$ épocas (implementar con PyTorch o JAX).
  \item Para cada arquitectura, graficar $f_{\text{MLP}}(t)$ junto
        con los datos y $\sin(2\pi t)$.
  \item Reportar $\mathcal{L}_{\text{train}}$ y
        $\mathcal{L}_{\text{test}}$ para cada una. Comparar con la
        tabla de la Sección~\ref{sec:oscilador4}.
  \item ¿Cuál arquitectura generaliza mejor? ¿Coincide con la
        predicción de la Proposición~\ref{prop:profundidad} (la red
        profunda (C) debería superar a (B) con parámetros comparables)?
\end{enumerate}

\medskip\noindent\textbf{Verificación:}
$(1, 8, 1)$: $\mathcal{L}_{\text{test}} \approx 0.06$ (subajuste leve:
pocas neuronas para capturar la forma completa).
$(1, 32, 1)$: $\mathcal{L}_{\text{test}} \approx 0.025$.
$(1, 16, 16, 1)$: $\mathcal{L}_{\text{test}} \approx 0.020$ (la red
profunda generaliza ligeramente mejor con más parámetros).

\medskip\noindent\textbf{Extensión:}
Repetir con activación tanh. ¿Cambia la calidad del ajuste? ¿Y la
velocidad de convergencia? Trazar las curvas de pérdida de
entrenamiento vs.\ época para ReLU y tanh.
\end{ejerciciocomp}

\begin{ejerciciocomp}{El MLP como aproximador de una función 2D}
\label{ejcomp:mlp_2d}
\textbf{Objetivo:} entrenar un MLP para aproximar una función de dos
variables con estructura no separable.

\medskip\noindent\textbf{Función objetivo:}
\begin{equation}
  g(u_1, u_2) = \exp\!\bigl(-0.5(u_1^2+u_2^2)\bigr)\,
  \cos(2\pi u_1)\cos(2\pi u_2),
  \quad (u_1,u_2) \in [-2,2]^2.
  \label{eq:funcion_2d}
\end{equation}
Esta función tiene oscilaciones locales moduladas por una envolvente
gaussiana: un campo con estructura multiescala.

\medskip\noindent\textbf{Algoritmo:}
\begin{enumerate}
  \item Generar $N_{\text{train}} = 500$ puntos uniformes en $[-2,2]^2$
        y $N_{\text{test}} = 1000$ puntos adicionales.
  \item Entrenar tres arquitecturas con activación tanh:
        (A)~$(2, 8, 1)$, (B)~$(2, 32, 1)$, (C)~$(2, 16, 16, 1)$.
        Usar MSE y Adam con $\eta = 10^{-3}$, 2000 épocas.
  \item Graficar $f_{\text{MLP}}(u_1, u_2)$ como mapa de calor para
        cada arquitectura. Comparar con $g$.
  \item Reportar $\mathcal{L}_{\text{test}}$ para cada una. ¿Cuál
        captura mejor la estructura multiescala?
\end{enumerate}

\medskip\noindent\textbf{Verificación:}
$(2, 8, 1)$: $\mathcal{L}_{\text{test}} \approx 0.04$ (no captura las
oscilaciones).
$(2, 32, 1)$: $\mathcal{L}_{\text{test}} \approx 0.008$.
$(2, 16, 16, 1)$: $\mathcal{L}_{\text{test}} \approx 0.003$ (la red
profunda es significativamente mejor para la estructura multiescala).
\end{ejerciciocomp}

% ─────────────────────────────────────────────────────────────────────────────
\begin{notasbib}
La neurona artificial como modelo del comportamiento neuronal fue
propuesta por McCulloch y Pitts (1943)~\cite{mcculloch1943logical}. El
perceptrón y su regla de aprendizaje fueron introducidos por Rosenblatt
(1958)~\cite{rosenblatt1958perceptron}. El teorema de aproximación
universal fue demostrado por Cybenko
(1989)~\cite{cybenko1989approximation} para activaciones sigmoide y
generalizado por Hornik, Stinchcombe y White
(1989)~\cite{hornik1989multilayer} para cualquier activación continua no
constante. Hornik (1991) mostró que la condición necesaria y suficiente
es la no polinomialidad de la activación.

El problema del gradiente que se desvanece fue identificado por
Hochreiter (1991) y analizado por Bengio, Simard y Frasconi (1994). La
activación ReLU como solución fue popularizada por Nair y Hinton
(2010)~\cite{nair2010rectified}. La revolución del aprendizaje profundo
con ReLU se asocia a la victoria de AlexNet (Krizhevsky, Sutskever e
Hinton, 2012)~\cite{krizhevsky2012imagenet} en ImageNet.

La conexión entre la sigmoide y la distribución de Fermi-Dirac, y entre
la tanh y la función de Langevin, es conocida en la literatura de
máquinas de Boltzmann~\cite{hinton2006fast}, donde la sigmoide emerge de
las probabilidades condicionales del modelo generativo. La cota de Barron
(1993) sobre la eficiencia de las superposiciones de sigmoides escapa la
maldición de la dimensionalidad: el error de aproximación es $O(1/m)$
independientemente de $n$, a diferencia de $O(m^{-2/n})$ para bases
polinomiales.

La ventaja de la profundidad sobre la anchura fue establecida
rigurosamente por Montúfar et al.\ (2014) para el conteo de regiones
lineales de redes ReLU, y por Telgarsky (2016) para la expresividad
composicional.
\end{notasbib}

% ─────────────────────────────────────────────────────────────────────────────
% ESTRELLA POLAR
% ─────────────────────────────────────────────────────────────────────────────

\begin{estrella}{El MLP como sustituto del flujo neutrónico}
\label{sec:estrella4}

\textbf{Del modelo lineal al MLP.}
La Estrella~\ref{sec:estrella2} mostró que la expansión modal lineal
$\phi(\bm{r};\bm{u}) \approx \sum_k a_k(\bm{u})\phi_k(\bm{r})$ con
$a_k$ lineal en $\bm{u}$ falla para configuraciones no lineales del
reactor. La alternativa es reemplazar los coeficientes lineales por un
MLP:
\[
  a_k(\bm{u}) = f_{\text{MLP}}^{(k)}(\bm{u};\bm{\theta}_k),
\]
donde $f_{\text{MLP}}^{(k)}$ es un MLP separado (o una salida de un MLP
compartido con $K$ salidas) que mapea los parámetros operacionales
$\bm{u} \in \mathbb{R}^5$ a los coeficientes modales.

\textbf{Arquitectura.}
Un MLP $(5, 64, 64, K)$ con $K = 20$ modos y activación tanh tiene
$p = (5+1)\cdot 64 + (64+1)\cdot 64 + (64+1)\cdot 20 = 384 + 4160 +
1300 = 5844$ parámetros. Con $N = 200$ simulaciones de referencia, la
relación $p/N \approx 29$ está en el régimen sobreparametrizado ($p \gg
N$), lo que requiere regularización (weight decay
$\lambda = 10^{-4}$, Capítulo~\ref{ch:cap2}) y monitorización del
sobreajuste (early stopping, Capítulo~\ref{ch:cap8}).

\textbf{Resultado.}
El MLP regularizado reduce el error de predicción del flujo neutrónico
de $\sim 2\%$ (expansión lineal regularizada) a $\sim 0.3\%$ en las 50
configuraciones de validación. La mejora es particularmente notable para
configuraciones con barras asimétricas, donde la expansión lineal tenía
error $> 5\%$.

\textbf{Limitación.}
El MLP proporciona predicciones deterministas: un solo valor
$\hat{\phi}(\bm{r};\bm{u})$ sin cuantificación de incertidumbre. En
aplicaciones de seguridad nuclear, saber cuánto confiar en la predicción
es tan importante como la predicción misma. El Capítulo~\ref{ch:cap7}
introducirá las herramientas bayesianas para cuantificar la
incertidumbre del sustituto, y la Estrella correspondiente las aplicará
al reactor.
\end{estrella}

\printbibliography[heading=subbibliography]